\section{Related Work}
\label{sec:related}

\Cref{tab:results} summarizes selected results on subtype entailment.

\begin{table}[tbp]
\caption{Selected results on subtype entailment. Earlier non-structural rows concern a
covariant constructor; our decidability result allows arbitrary variance and arity.
The complexity bounds for our results are proved in \cref{sec:complexity};
\cref{sec:mechanization} states which parts of them are mechanized.}
\label{tab:results}
\small\setlength{\tabcolsep}{3pt}
\begin{tabular}{@{}>{\raggedright\arraybackslash}p{0.29\linewidth}>{\raggedright\arraybackslash}p{0.18\linewidth}>{\raggedright\arraybackslash}p{0.20\linewidth}>{\raggedright\arraybackslash}p{0.27\linewidth}@{}}
\toprule
Problem & Trees & Result & Reference \\
\midrule
Structural entailment & finite & coNP-complete & \citet{HengleinRehof97} \\
Structural entailment & arbitrary & PSPACE-complete & \citet{HengleinRehof98} \\
Structural, first-order theory & finite & decidable & \citet{KuncakRinard03} \\
\NSSE, $\ftree$ unary & finite, regular, arbitrary & decidable & \citet{NiehrenPriesnitz03} \\
\NSSE, $\ftree$ binary, no literals $x = \Bot$, $x = \Top$ & finite, arbitrary & PSPACE-complete & \citet{NiehrenPriesnitz99} \\
\NSSE, $\ftree$ binary & finite, arbitrary & PSPACE-hard & \citet[Theorems~1--2]{HengleinRehof98} \\
\NSSE, any arity and variance & finite, regular, arbitrary & decidable; in 2-EXPSPACE (stronger for $n\le1$, \cref{app:small-arity}) & this paper \\
\NSSE, $\ftree$ covariant, any fixed arity $n\ge2$ & finite, regular, arbitrary & PSPACE-hard & this paper \\
Non-structural, first-order theory & finite, regular, arbitrary & undecidable & \citet{SuAikenNiehrenPriesnitzTreinen02} \\
\bottomrule
\end{tabular}
\end{table}

\paragraph{Subtype Entailment and Constraint Simplification}
Subtype constraints arise in type inference with subtyping
\citep{Mitchell91,FuhMishra90,AikenWimmers93,EifrigSmithTrifonov95,PalsbergWandOKeefe97},
and the size of inferred constraint sets made simplification a central practical concern
\citep{Pottier96,Pottier98thesis,Pottier01,TrifonovSmith96}. Sound but incomplete
entailment tests based on syntactic closure are used in practice; \citet{Pottier01} discusses
their limits. Satisfiability of subtype
constraints is comparatively well understood
\citep{Tiuryn92,PrattTiuryn96,Frey97,KozenPalsbergSchwartzbach94,PalsbergOKeefe95,PalsbergWandOKeefe97,Simonet03,NiehrenPriesnitzSu05}.
Closure-based satisfiability characterizations are classical.
\citet{KozenPalsbergSchwartzbach94} treat a language with $\Top$ but
without $\Bot$; related characterizations appear in
\citet{PalsbergOKeefe95,PalsbergWandOKeefe97}.
Here the additional ingredients are a common least shape under variable flip
and polarity selection for variance. Subtyping between recursive types without variables
is also well understood
\citep{AmadioCardelli93,KozenPalsbergSchwartzbach95}. For entailment, \citet{HengleinRehof97} proved coNP-completeness for structural subtyping over
finite trees and \citet{HengleinRehof98} PSPACE-completeness over arbitrary trees. Their Theorems~1 and~2 establish
non-structural PSPACE-hardness using a binary covariant constructor, over
arbitrary and finite trees, respectively; this is the scope recorded in
\cref{tab:results}. \Cref{thm:cx-hard} gives a different reduction, from
universality of nondeterministic automata, with a uniform proof for all three
domains and every fixed arity $n\ge2$. \citet{Rehof97,Rehof98thesis} studies
minimal typings and the complexity of atomic and structural subtyping.
\citet{KuncakRinard03} show that the full first-order theory of structural subtyping of
non-recursive types is decidable.

\paragraph{Non-Structural Subtype Entailment}
\citet{NiehrenPriesnitz99} proved PSPACE-completeness for the fragment in which the
constraints contain no literals $x = \Bot$ and $x = \Top$ (the constants may still occur in
solutions), by a reduction to universality of finite automata; they introduced P-automata to
explain why the nonregular sets of paths that arise in general can be ignored in this
fragment. \citet{NiehrenPriesnitz01,NiehrenPriesnitz03} introduced cap
automata and cap set expressions, proved decidability of the unary case,
and identified the role of word equations in the general problem. Their
forward reduction from entailment to cap-automaton universality motivates
our use of periodic suffix comparisons. \Cref{sec:errata} discusses the
proof dependencies of their converse reduction and arity equivalence.
Priesnitz's theses \citep{Priesnitz00,Priesnitz04} develop the automata
approach further.
\citet{SuAikenNiehrenPriesnitzTreinen02} proved that the first-order theory
of non-structural subtyping is undecidable. This does not settle entailment,
which is a restricted universal implication. They and \citet{Su02} also
proposed a reduction of entailment to constrained tree automata on tuples of
trees, whose emptiness problem is undecidable in general \citep{Treinen00};
an extended version of the paper develops this reduction further
\citep{SuEtAl06}. A preprint
claimed membership in PSPACE for the finite and regular cases \citep{Schubert21}; it was
withdrawn by its author because of an error in a key lemma. The
problem has been described as open in work on intersection types and on type inference
\citep{DudenhefnerMartensRehof17,JiangCuiOliveira25}.

\subsection{The Published Characterization}
\label{sec:errata}
Proposition~16 of NP03 is false as stated:
\cref{app:np03} gives a six-state complete restricted cap automaton that
rejects $11$, whereas the left automaton of its back-translated judgment
accepts $11$ under either finality convention for Table~3. The judgment
is moreover a satisfiable entailment. Theorem~4 and the converse
reduction underlying Theorem~1 rely on Proposition~16.
Corollary~1 (equivalence of arities at least two and of tree classes)
also relies on that converse, and its arity part on Proposition~17.
Theorem~4, Theorem~1, and Corollary~1 themselves are not refuted;
only their published proofs are incomplete. No defect in Proposition~17
is claimed here. The forward characterization
(Theorem~2, with Propositions~8 and~11) is unaffected: it is
all that NP03's unary decidability result and the earlier binary proof
in \cref{sec:earlier} use.
The proofs of Propositions~12 and~14 treat states ready by the bottom,
top, or reflexivity rule as final. Table~3 must be read with that
convention: read literally, it rejects the path $1$ of NP03's unary
example $\judg{z=\Top\wedge f(z)\le y}{x}{y}$ (Section~4), which NP03
shows to be both left safe and right safe. \Cref{app:np03} gives the
acceptance details.
Our present proof uses neither the forward characterization nor any of
these converse or alphabet reductions.

\subsection{Other Connections}
\paragraph{Set Constraints and Feature Trees}
Entailment has been studied for neighboring constraint languages. Entailment of atomic set
constraints is PSPACE-complete over infinite signatures \citep{NiehrenMullerTalbot99}.
Ordering constraints over feature trees admit efficient entailment tests under an
assumption of infinitely many features \citep{MullerNiehrenPodelski00}, while their full
first-order theory is undecidable \citep{MullerNiehrenTreinen98}. Componential set-based
analysis uses entailment of set constraints for simplification \citep{FlanaganFelleisen99}.

\paragraph{Type Inference Without General Entailment}
Several modern approaches to type inference with subtyping avoid general entailment by
changing the type language or the semantics. Algebraic subtyping
\citep{Dolan16,DolanMycroft17,Parreaux20} interprets types in an algebraic model rather than
as ground trees, and infers principal types by biunification on polarized types, treating
subsumption between inferred types separately. MLstruct extends this approach to a Boolean
algebra of structural types \citep{ParreauxChau22}, and \citet{ParreauxEtAl24} retain
subtyping constraints in polymorphic types. Semantic subtyping interprets types as sets of
values and decides subtyping between types, including with polymorphism
\citep{CastagnaXu11,GesbertGenevesLayaida15}. Orthologic type systems obtain efficient
entailment for an open-world semantics \citep{GuilloudKuncak25}. These approaches use different type languages or semantic models, so their
decidability results do not directly decide the ground-tree entailment
problem studied here.

\paragraph{Word Equations and Periodicity}
Satisfiability of word equations is decidable \citep{Makanin77} and in PSPACE
\citep{Plandowski04}, but the positive $\forall\exists^3$ fragment is undecidable
\citep{Durnev95}. \citet[Corollary~2]{NiehrenPriesnitz03} express \NSSE in the positive
$\forall\exists^*$ fragment with regular constraints. Existential word-equation algorithms do not
directly decide this quantified formulation. Undecidability of the ambient
fragment leaves open the restricted instances arising from \NSSE, which
our reader construction and compression argument decide.
The periodicity arguments use the principle that sufficiently long agreement
between two periods forces agreement between their residue classes
\citep{FineWilf65}. \Cref{lem:rotation} proves the special case it needs
directly, and in \cref{lem:truncation}, equality over one full period
propagates along the periodic tail. These are instances of classical
periodic-word reasoning \citep{Lothaire02}, without a general word-equation
algorithm.

\paragraph{Back-References}
A cap repeats a root word that is chosen during the run, which resembles a back-reference in
an extended regular expression. For extended regular expressions, universality is
undecidable even with a single variable \citep{Freydenberger11}. Cap set expressions are a
much weaker mechanism: the only repetition is the prefix closure of the powers of one word,
and that word must be followed to the end of the path. Even so, universality of unrestricted
cap automata was left open by \citet[Section~10]{NiehrenPriesnitz03};
our decision procedure applies to readers satisfying the rejected-path property
and does not settle that question.

\paragraph{Desubstitution}
Cutting a sequence at the occurrences of a word and recoding the blocks between them is the
idea behind \emph{return words} and \emph{derived sequences} in symbolic dynamics
\citep{Durand98}, and the question whether such decompositions are unique is the subject of
recognizability theory \citep{Mosse92}. Our hierarchy cuts a finite word at its last letter,
and \cref{lem:cores} identifies the retained positions with the occurrences of a terminal
word. The connection is the recoding of intervals between repeated
occurrences. Here the chosen marker directly delimits the blocks, so unique
parsing needs no recognizability theorem. We use the hierarchy as a measure of the
complexity of a hole whose levels can be bounded (\cref{thm:depth}) and compressed
(\cref{thm:search}).

\paragraph{Automata with Constraints}
Our problem resembles emptiness for automata with equality and disequality constraints, where
pumping must preserve the constraints. For tree automata with constraints between brothers
\citep{BogaertTison92}, reduction automata \citep{DauchetCaronCoquide95}, and global constraints
\citep{TATA,ComonJacquemard03,FiliotTalbotTison10,BargunoEtAl10}, emptiness proofs replace
subterms \emph{simultaneously} while preserving states and equality patterns, rather than
pumping one subterm at a time. The column summaries of \cref{sec:search} follow the same
principle for the blocks of one hierarchy level. The new difficulty in \NSSE is that the
constraints compare suffixes of a single word at unbounded offsets; the admission cone
and the depth bound are what reduce it to a setting where simultaneous replacement suffices.
